\documentclass[11pt,a4paper]{article}

% ─── Pacotes ────────────────────────────────────────────────────────────────
\usepackage[utf8]{inputenc}
\usepackage[T1]{fontenc}
\usepackage[top=3.5cm, bottom=2.5cm, left=2.5cm, right=2.5cm]{geometry}
\usepackage{xcolor}
\usepackage{titlesec}
\usepackage{enumitem}
\usepackage{amsmath}
\usepackage{mdframed}
\usepackage{graphicx}
\usepackage{fancyhdr}
\usepackage{tcolorbox}
\usepackage{tabularx}
\usepackage{booktabs}
\usepackage{microtype}
\usepackage{hyperref}

\tcbuselibrary{skins, breakable}

% ─── Cores ──────────────────────────────────────────────────────────────────
\definecolor{quantpurple}{HTML}{7B2FBE}
\definecolor{quantpurplelight}{HTML}{F0E8FF}
\definecolor{quantdark}{HTML}{1A1A1A}
\definecolor{quantgray}{HTML}{6B7280}
\definecolor{quantlightgray}{HTML}{F3F4F6}
\definecolor{quantgreen}{HTML}{16A34A}
\definecolor{quantred}{HTML}{DC2626}
\definecolor{quantamber}{HTML}{D97706}

% ─── Fonte sem serifa ────────────────────────────────────────────────────────
\renewcommand{\familydefault}{\sfdefault}

% ─── Hyperlinks ─────────────────────────────────────────────────────────────
\hypersetup{hidelinks}

% ─── Cabeçalho e rodapé ─────────────────────────────────────────────────────
\setlength{\headheight}{2cm}
\setlength{\headsep}{1cm}

\pagestyle{fancy}
\fancyhf{}
\renewcommand{\headrulewidth}{0pt}
\fancyhead[L]{%
  \includegraphics[height=1.40cm]{logoSimples.png}
}
\fancyhead[R]{%
  \raisebox{0.5cm}{\textcolor{quantgray}{\footnotesize Núcleo Quant
  \,$\cdot$\, Liga de Mercado Financeiro
  \,$\cdot$\, UFU \,$\cdot$\, 2026.1}}
}
\fancyfoot[C]{\textcolor{quantgray}{\footnotesize\thepage}}
\fancyfoot[R]{%
  \textcolor{quantgray}{\footnotesize Aula S2 \,$\cdot$\, Análise Exploratória de Ativos}
}

% ─── Títulos de seção ───────────────────────────────────────────────────────
\titleformat{\section}
  {\large\bfseries\color{quantpurple}}
  {}{0em}{}
  [\vspace{2pt}\textcolor{quantpurple}{\rule{\textwidth}{1.2pt}}\vspace{4pt}]

\titlespacing{\section}{0pt}{18pt}{6pt}

% ─── Listas ─────────────────────────────────────────────────────────────────
\setlist[itemize]{
  label=\textcolor{quantpurple}{\small$\blacktriangleright$},
  leftmargin=1.8em,
  itemsep=3pt,
  parsep=0pt
}
\setlist[enumerate]{
  label=\textcolor{quantpurple}{\textbf{\arabic*.}},
  leftmargin=1.8em,
  itemsep=5pt,
  parsep=0pt
}

% ─── Código inline ──────────────────────────────────────────────────────────
\newcommand{\code}[1]{%
  \colorbox{quantlightgray}{\texttt{\small\color{quantdark}#1}}%
}

% ─── Caixa de destaque (roxo claro) ─────────────────────────────────────────
\newtcolorbox{highlightbox}{
  enhanced,
  colback=quantpurplelight,
  colframe=quantpurple,
  boxrule=1pt,
  arc=4pt,
  left=10pt, right=10pt, top=7pt, bottom=7pt,
  fontupper=\small\color{quantdark}
}

% ─── Caixa de passo ─────────────────────────────────────────────────────────
\newtcolorbox{stepbox}[2][]{
  enhanced,
  colback=white,
  colframe=quantpurple,
  boxrule=0.8pt,
  arc=4pt,
  title={\textbf{\textcolor{white}{#2}}},
  attach boxed title to top left={yshift=-1mm, xshift=10pt},
  boxed title style={
    colback=quantpurple,
    colframe=quantpurple,
    arc=3pt,
    boxrule=0pt,
    left=6pt, right=6pt, top=2pt, bottom=2pt
  },
  left=10pt, right=10pt, top=14pt, bottom=8pt,
  #1
}

% ─── Caixa de atenção (âmbar) ───────────────────────────────────────────────
\newtcolorbox{warningbox}{
  enhanced,
  colback=white,
  colframe=quantamber,
  boxrule=0.8pt,
  arc=4pt,
  left=10pt, right=10pt, top=6pt, bottom=6pt,
  fontupper=\small\color{quantdark}
}

% ════════════════════════════════════════════════════════════════════════════
\begin{document}

% ════════════════════════════════════════════════════════════════════════════

\thispagestyle{fancy}

% ─── Título ─────────────────────────────────────────────────────────────────
\begin{center}
  \includegraphics[height=3cm]{logoLiga.png}\\[1em]
  {\fontsize{22}{26}\selectfont\textbf{EXERCÍCIO PRÁTICO}}\\[0.4em]
  {\large\textcolor{quantgray}{Análise Exploratória de Ativos com Python}}
\end{center}

\vspace{0.5em}
\textcolor{quantpurple}{\rule{\textwidth}{2pt}}
\vspace{1em}

% ─── Contexto ───────────────────────────────────────────────────────────────
\section{Contexto}

Um cliente solicita uma análise quantitativa de quatro ativos de sua carteira ao longo dos
últimos quatro anos. O objetivo é ir além da simples observação de preços e construir uma
visão estruturada sobre retorno, risco e correlação entre os ativos.

\begin{highlightbox}
  \textbf{Premissa fundamental:} preços absolutos não são comparáveis entre ativos distintos.
  Uma ação de R\$\,200 não é necessariamente mais cara ou mais rentável do que uma de R\$\,30.
  A análise correta parte de \textbf{retornos percentuais}, que nivelam todos os ativos em uma
  escala comum e permitem comparações válidas.
\end{highlightbox}

\vspace{0.8em}

\noindent\textbf{Ativos a serem analisados:}

\vspace{0.5em}
\begin{tabularx}{\textwidth}{@{}llX@{}}
  \toprule
  \textcolor{quantpurple}{\textbf{Ticker}} &
  \textcolor{quantpurple}{\textbf{Ativo}} &
  \textcolor{quantpurple}{\textbf{Descrição}} \\
  \midrule
  \code{MGLU3.SA}  & Magazine Luiza  & Setor de varejo; histórico de alta volatilidade no período \\[3pt]
  \code{WEGE3.SA}  & WEG S.A.        & Setor industrial; trajetória de crescimento consistente \\[3pt]
  \code{BOVA11.SA} & ETF Ibovespa    & Benchmark do mercado acionário brasileiro \\[3pt]
  \code{USDBRL=X}  & Dólar (USD/BRL) & Ativo de proteção cambial; hedge natural em crises \\
  \bottomrule
\end{tabularx}

\vspace{0.8em}
\noindent\textbf{Período de análise:}\quad
\code{2020-01-01} \textcolor{quantgray}{a} \code{2024-01-01}

\vspace{1.2em}

\begin{warningbox}
  \textbf{\textcolor{quantamber}{Atenção ao baixar o Dólar:}} o ticker \code{USDBRL=X} opera em
  calendário forex, diferente do calendário de pregão da B3. Realize o download \textbf{separadamente}
  e utilize \code{join(how="left")} para alinhar os dados ao calendário brasileiro. Consulte o
  Passo~1 para o procedimento correto.
\end{warningbox}

\newpage

% ─── Etapas ──────────────────────────────────────────────────────────────────
\section{Etapas}

% ── PASSO 1 ──────────────────────────────────────────────────────────────────
\begin{stepbox}{Passo 1 \quad Coleta e Tratamento de Dados}
\begin{enumerate}
  \item Utilize \code{yfinance} para baixar os dados de fechamento ajustado (\code{Close}) dos
        quatro ativos. O parâmetro \code{auto\_adjust=True} deve ser utilizado para que splits e
        pagamentos de dividendos não gerem quedas artificiais no histórico de preços.

  \item O ticker \code{USDBRL=X} deve ser baixado \textbf{separadamente}. Antes de unir os
        DataFrames, remova a informação de fuso horário dos índices com \code{.tz\_localize(None)}
        e utilize \code{join(how="left")} para preservar apenas os dias de pregão da B3.

  \item Aplique o tratamento de dados faltantes em duas etapas:
        \begin{itemize}
          \item \code{ffill()} para propagar o último preço válido em feriados e dias sem negociação
          \item \code{dropna()} para remover as linhas iniciais sem histórico disponível
        \end{itemize}
\end{enumerate}
\end{stepbox}

\vspace{1em}

% ── PASSO 2 ──────────────────────────────────────────────────────────────────
\begin{stepbox}{Passo 2 \quad Normalização e Retornos}
\begin{enumerate}
  \item Normalize os preços para \textbf{Base 100} e gere um gráfico de linhas interativo com
        Plotly. A normalização permite comparar ativos em escalas diferentes no mesmo eixo:

        \[
          \text{df\_norm} = \frac{\text{df\_clean}}{\text{df\_clean.iloc[0]}} \times 100
        \]

  \item Calcule os retornos simples diários utilizando \code{pct\_change()}:

        \[
          R_t = \frac{P_t - P_{t-1}}{P_{t-1}}
        \]

        Retornos percentuais são a métrica adequada para comparação entre ativos e para
        a construção de qualquer análise estatística subsequente.
\end{enumerate}
\end{stepbox}

\vspace{1em}

% ── PASSO 3 ──────────────────────────────────────────────────────────────────
\begin{stepbox}{Passo 3 \quad Análise de Volatilidade}
\begin{enumerate}
  \item Calcule a \textbf{volatilidade anualizada} de cada ativo, definida como o desvio-padrão
        dos retornos diários multiplicado pela raiz quadrada de 252 dias úteis:

        \[
          \sigma_{\text{anual}} = \sigma_{\text{diária}} \times \sqrt{252}
        \]

        Identifique qual ativo apresenta maior e menor volatilidade no período. Interprete
        o que esses valores implicam do ponto de vista do risco assumido pelo investidor.

  \item Gere um \textbf{histograma interativo} com os retornos diários de \code{MGLU3.SA} e
        \code{WEGE3.SA} sobrepostos. Analise a forma das distribuições: ativos com maior risco
        apresentam caudas mais pesadas e distribuição mais dispersa em torno da média.
\end{enumerate}
\end{stepbox}

\newpage

% ── PASSO 4 ──────────────────────────────────────────────────────────────────
\begin{stepbox}{Passo 4 \quad Correlação e Diversificação}
\begin{enumerate}
  \item Calcule a \textbf{matriz de correlação} dos retornos simples com \code{.corr()}.
        Correlações próximas de $+1{,}0$ indicam que os ativos se movem de forma conjunta,
        reduzindo o efeito da diversificação em períodos de estresse.

  \item Gere um \textbf{heatmap interativo} com Plotly Express para visualizar as correlações:

        \begin{highlightbox}
          \small\ttfamily
          fig = px.imshow(corr, text\_auto=".2f",\\
          \hspace{4.4em} color\_continuous\_scale="RdBu\_r",\\
          \hspace{4.4em} zmin=-1, zmax=1)
        \end{highlightbox}

        Analise a relação entre o Dólar e o Ibovespa. Com base no que foi discutido em aula,
        qual o sinal esperado para essa correlação e qual a sua implicação para a gestão de risco?
\end{enumerate}
\end{stepbox}

\vspace{1em}

% ── PASSO 5 ──────────────────────────────────────────────────────────────────
\begin{stepbox}{Passo 5 \quad Médias Móveis e Análise de Tendência}
\begin{enumerate}
  \item Isole os dados de fechamento de \code{WEGE3.SA} e calcule as médias móveis simples:
        \begin{itemize}
          \item \textbf{MM20} (20 dias): \code{close.rolling(20).mean()}
          \item \textbf{MM200} (200 dias): \code{close.rolling(200).mean()}
        \end{itemize}

  \item Plote as três séries no mesmo gráfico interativo e identifique pelo menos um dos eventos:
        \begin{itemize}
          \item \textbf{\textcolor{quantgreen}{Golden Cross:}} MM20 cruza a MM200 para cima,
                indicando possível início de tendência de alta
          \item \textbf{\textcolor{quantred}{Death Cross:}} MM20 cruza a MM200 para baixo,
                indicando possível início de tendência de baixa
        \end{itemize}

        Utilize \code{add\_vline()} no Plotly para sinalizar os cruzamentos detectados.
        Discuta brevemente as limitações desse tipo de sinal antes de utilizá-lo em uma estratégia.
\end{enumerate}
\end{stepbox}

\vspace{1.8em}

% ─── Rodapé do documento ─────────────────────────────────────────────────────
\textcolor{quantpurple}{\rule{\textwidth}{1.2pt}}

\vspace{0.6em}
\begin{center}
  \textcolor{quantgray}{\small
    Entrega via GitHub \,$\cdot$\, repositório com \texttt{README} descrevendo o que foi
    implementado e instruções para execução \,$\cdot$\, individual ou em dupla
  }\\[5pt]
  \textcolor{quantgray}{\small
    Critérios de avaliação:\;
    \textbf{Rigor Técnico} \,$\cdot$\,
    \textbf{Clareza} \,$\cdot$\,
    \textbf{Reprodutibilidade} \,$\cdot$\,
    \textbf{Interpretação dos Resultados}
  }
\end{center}

\end{document}